\section{Lambert expansions and explicit
coefficients}\label{tcl:lambert-expansions-and-explicit-coefficients}

All limits are for positive integer \(n\) tending to infinity. The
polynomial \(P\) and expansion order \(J\) are fixed before taking that
limit. Implicit constants can depend on \(P\) and \(J\). Numerical
diagnostics below are not rigorous intervals.


\subsection{Definitions and main
theorem}\label{tcl:definitions-and-main-theorem}

Write \((u)_k=u(u-1)\cdots(u-k+1)\) for a falling factorial and
\(u^{\overline{k}}=u(u+1)\cdots(u+k-1)\) for a rising factorial. Put

\[
\begin{gathered}
 M(m)=m(m-1)/2,\qquad
 H_n=e^{-1}\sum_{m\ge0}\frac{(M(m))_n}{m!},\\
 c_k=[z^k]e^{P(z)},\qquad
 A_P(n)=\sum_{k=0}^n(n)_kc_kH_{n-k}.
\end{gathered}
\]

Assume \(P\) is a fixed real polynomial with \(P(0)=0\). The relevant
cases are

\[
 P_V(z)=-z/2,\quad P_U(z)=z/2,\quad
 P_L(z)=z/2-z^3/6-z^4/4-z^5/8-z^6/48.
\]

Let

\[
 r=W(2n),\qquad s=e^r=2n/r,
\]

\[
 S_n=\frac{\exp\{2n(r-1+1/r)-1-n\log2-r^2/4-r/2\}}{\sqrt{1+r}}.
 \tag{C1}\label{tcl:eq:C1}
\]

\textbf{Theorem C1.} If \(K_P\) is the field generated over the rationals
by the coefficients of \(P\), there are rational functions
\(a_{P,j}(r)\) in \(K_P(r)\), obtained by the finite algebraic
prescription below, with \(a_{P,0}=1\) and

\[
 a_{P,j}(r)=O_{P,j}((1+r)^{4j}),
\]

such that, for every fixed \(J\ge0\),

\[
 A_P(n)=S_n\left(\sum_{j=0}^J\frac{a_{P,j}(r)}{n^j}
 +O_{P,J}\!\left(\frac{(1+r)^{4J+4}}{n^{J+1}}\right)\right).
 \tag{C2}\label{tcl:eq:C2}
\]

For rational \(P\) these are ordinary rational functions over the
rationals and the prescription is exactly computable by rational
arithmetic. For arbitrary real \(P\) the prescription is algebraic
relative to its coefficients; no computable access to an arbitrary real
constant is assumed.

Thus the statement covers every fixed order. The accompanying software
has been run and checked through \(J=3\); computation at three orders is
not the proof of the all-orders assertion. The public interface supports
orders zero through three. The finite algebraic prescription itself has
no fixed mathematical order cap.

\subsection{Exact finite-order coefficient
generator}\label{tcl:exact-finite-order-coefficient-generator}

Use formal \(h\) and \(y\), with

\[
 h=s^{-1/2},\quad x=1+hy,\quad n=r/(2h^2),\quad
 m=x/h^2,\quad M=\frac{x^2}{2h^4}(1-h^2/x).
\]

Write $\mathsf B_k(z)$ for the Bernoulli polynomial and
$\mathsf B_k=\mathsf B_k(0)$ for the Bernoulli number. These differ from
the Bell numbers $B_k$ used earlier; for example $\mathsf B_2=1/6$.
For \(p\ge1\) let

\[
 F_p(n)=\sum_{j=0}^{n-1}j^p
       =\frac{\mathsf B_{p+1}(n)-\mathsf B_{p+1}(0)}{p+1}.
\]

The identity in the second expression defines its polynomial
continuation. Set

\[
 D=n\log(1-h^2/x)-\sum_{p\ge1}\frac{F_p(n)}{pM^p},
 \tag{C3}\label{tcl:eq:C3}
\]

\[
 E=\frac{(r-x)\log x+(1-r)(x-1)}{h^2}
 +\frac{r+1}{2}y^2-\frac12\log x
 -\sum_{q\ge1}\frac{\mathsf B_{2q}h^{4q-2}}{2q(2q-1)x^{2q-1}}
 +D+\frac r2+\frac{r^2}4,
 \tag{C4}\label{tcl:eq:C4}
\]

\[
 Q=\sum_{k\ge0}c_k\frac{(n)_k}{(M-n+1)^{\overline{k}}}.
 \tag{C5}\label{tcl:eq:C5}
\]

Equations (C3)--(C5) are finite-order formal instructions. To obtain
\(a_{P,j}\), keep \(p\le j+1\) in (C3), \(q\) with \(4q-2\le2j\) in
(C4), and \(k\le j\) in (C5), and discard powers above \(h^{2j}\). In
particular, no infinite Stirling series is asserted to converge.

If \(Y\) is centered Gaussian with variance \(1/(r+1)\), then

\[
 a_{P,j}(r)=\left(\frac r2\right)^j
 \mathbb E\left([h^{2j}]e^{E(h,Y,r)}Q(h,Y,r)\right).
 \tag{C6}\label{tcl:eq:C6}
\]

This is an exact algebraic prescription over \(K_P\), and is an exact
rational-arithmetic prescription when \(P\) has rational coefficients.
Gaussian moments are

\[
 \mathbb E Y^{2\ell}=(2\ell-1)!!/(1+r)^\ell,
 \qquad \mathbb E Y^{2\ell+1}=0.
\]

An efficient implementation computes the coefficients of \(\exp(E)\) by
the finite recurrence

\[
 b_0=1,\qquad b_j=\frac1j\sum_{k=1}^j kE_kb_{j-k},
 \quad E=\sum_{k\ge1}E_kh^k.
\]

The code \texttt{generate\_coefficients.py} implements this using exact
SymPy rational arithmetic for the displayed rational polynomials \(P\);
formal polynomial parameters can instead be retained symbolically.
Setting \(D=0\) and omitting its centering constants gives the usual
Bell saddle corrections \(b_j\) for

\[
 2^{-n}B_{2n}=\frac{e^{2n(r-1+1/r)-1-n\log2}}{\sqrt{1+r}}
 \left(\sum_{j=0}^J b_j(r)n^{-j}
 +O_J((1+r)^{4J+4}n^{-J-1})\right).
\]

The coefficients relative to the exact Bell factor are obtained by
ordinary formal division, not by identifying \(S_n\) with that exact
factor.

\subsection{Justification of the fixed-order
remainder}\label{tcl:justification-of-the-fixed-order-remainder}

\subsubsection{Localization and signed
convolution}\label{tcl:localization-and-signed-convolution}

The preceding exact-saddle section establishes the following facts
directly from strict concavity of

\[
 f_n(t)=\sum_{j<n}\log(M(t)-j)-\log\Gamma(t+1),
 \qquad t>(1+\sqrt{1+8(n-1)})/2.
\]

Its unique saddle \(\tau\) satisfies \(\tau=s+O(r)\), its variance is
\(\sigma^2=s(1+O(r^2/n))/(r+1)\), and on an interval comparable with
\(s\) the exponent loses at least \(c(t-\tau)^2/\sigma^2\). Outside that
interval a concave tangent gives an exponentially decreasing right tail;
the left tail has only \(O(s)\) integer points. In particular
\(H_n/S_n=1+O(r^4/n)\). The same estimates hold with \(n\) replaced by
\(n-k\) for any fixed \(k\).

Positivity also gives the adjacent-ratio upper bound

\[
 H_{n-1}/H_n\le C(\log n/n)^2.
\]

Combining it for \(k\le n/2\) with the entire majorant
\(\exp(\sum_j|[z^j]P|z^j)\), and using a coefficient Cauchy bound for
\(k>n/2\), yields

\[
 A_P(n)=\sum_{k=0}^J(n)_kc_kH_{n-k}
 +O_{P,J}(H_n(r^2/n)^{J+1}).
 \tag{C7}\label{tcl:eq:C7}
\]

This is an absolute bound and does not assume the \(c_k\) have one sign.
For \(d=\deg P\ge1\) the normalized absolute far tail \(k>n/2\) is at
most

\[
 \exp\{-\tfrac12(1+1/d)n\log n+n\log\log n+O_P(n)\}.
\]

Therefore only fixed \(k\le J\) are needed for the Lambert coefficient
calculation.

For each such \(k\), discard its extra left support \(n-k\le M(m)<n\).
There \(m=O(\sqrt n)\), and the logarithm of its unnormalized summand is
at most \(n\log n+O_J(n+\log n)\), whereas
\(\log H_n=2n\log n-2n\log\log n+O(n)\). It is exponentially negligible
relative to every fixed algebraic error required here. On the common
support \(M(m)\ge n\), the exact identity

\[
 \frac{(M)_{n-k}}{(M)_n}
  =\frac1{(M-n+1)^{\overline{k}}}
 \tag{C8}\label{tcl:eq:C8}
\]

justifies the finite version \(Q_J\) of (C5). Applying (C8) to the
discarded support would be invalid.

Nor is a uniform bound on \(Q_J\) near \(M=n\) needed. Localize each
original positive sum \(H_{n-k}\) separately before introducing \(Q_J\).
For fixed \(k\) its saddle differs from \(s\) by \(O_J(r)\), and its
variance is comparable with \(s/(1+r)\). Thus the portion outside the
common Lambert window \(|Z|\le r^2\) is \(O_J(H_{n-k}e^{-c_Jr^4})\).
Multiplying by \((n)_k|c_k|\) and using
\((n)_kH_{n-k}/H_n=O_J((r^2/n)^k)\) gives \(O_{P,J}(H_ne^{-c_Jr^4})\)
after summing \(k\le J\). This controls the boundary and far tails even
when the rational multiplier itself is large there. Only inside the
common central window is \(Q_J\) expanded analytically.

\subsubsection{A finite analytic
exponent}\label{tcl:a-finite-analytic-exponent}

Fix \(J\). In (C3), retain \(p=1,\ldots,J+1\) and retain the logarithm
\(n\log(1-h^2/x)\) analytically, with its removable value at \(h=0\). In
(C4), retain only those Stirling terms with \(4q-2\le2J\). Call the
resulting finite analytic exponent \(E_J\). Use \(Q_J\) from (C5) with
\(k\le J\).

The original real collision logarithm differs from this finite version,
uniformly for \(m\) comparable with \(s\), by

\[
 O_J(n^{J+3}/M^{J+2})
 =O_J(r^{J+3}h^{2J+2}).
 \tag{C9}\label{tcl:eq:C9}
\]

Indeed the omitted logarithm tail is bounded by a constant times
\(\sum_{j<n}(j/M)^{J+2}\), since \(n/M=o(1)\). The omitted real
positive-argument Stirling remainder is bounded by its first omitted
term, and hence is \(O_J(h^{2J+2})\) or smaller. These are estimates for
real \(h\); they are not deduced from convergence of the Stirling
series. The standard finite Stirling remainder is documented in
\href{https://dlmf.nist.gov/5.11}{DLMF §5.11}.

At \(h=0\), \(D\) begins with \(-r/(2x)-r^2/(4x^2)\), so
\(E_J(0,y,r)=0\). Every singular-looking power of \(h\) in the finite
expression is removable. For example, \(F_p(n)/M^p\) starts at
\(h^{2p-2}\); the quadratic term of the first expression in (C4) cancels
its explicitly added Gaussian term.

\subsubsection{Cauchy radius, integrated remainder, and
parity}\label{tcl:cauchy-radius-integrated-remainder-and-parity}

Rescale \(Z=\sqrt{1+r}\,y\). For real \(Z\) and \(r\ge1\), there is a
sufficiently small \(c_J>0\) such that on the complex \(h\)-disk

\[
 |h|\le R_Z=c_J(1+r)^{-3/2}(1+|Z|)^{-3}
 \tag{C10}\label{tcl:eq:C10}
\]

we have

\[
 |e^{E_J(h,Z/\sqrt{1+r},r)}Q_J(h,Z/\sqrt{1+r},r)|\le C_{P,J}.
 \tag{C11}\label{tcl:eq:C11}
\]

Here is a direct bound behind (C11). On this disk,
\(|hy|\le c_J(1+r)^{-2}|Z|/(1+|Z|)^3\), so all denominators \(x\) and
the finitely many factors in \(Q_J\) stay away from zero. The shifted
leading collision part is \(O(r^2|hy|)\), hence bounded. The next
collision contribution is \(O(r^3|h|^2)\), hence bounded, and all
subsequent fixed collision terms are smaller. The cubic-and-higher Bell
phase is \(O((1+r)|h||y|^3/(1-|hy|))\), also bounded. The logarithmic
amplitude and the finite Stirling terms are bounded as well. The finite
\(Q_J\) is bounded since each \(k\)-term starts at \(O((r|h|^2)^k)\) and
all its normalized denominators are nonzero. Constants may be reduced or
enlarged for finitely many lower powers and fixed coefficients.

Cauchy's estimate therefore bounds its coefficient of \(h^q\) by

\[
 C_{P,J,q}(1+r)^{3q/2}(1+|Z|)^{3q}.
 \tag{C12}\label{tcl:eq:C12}
\]

On the real central window \(|Z|\le r^2\), the actual \(h=e^{-r/2}\)
eventually satisfies \(h\le R_Z/2\). Expanding through degree \(2J+1\)
has pointwise remainder bounded by

\[
 C_{P,J}h^{2J+2}(1+r)^{3J+3}(1+|Z|)^{6J+6}.
 \tag{C13}\label{tcl:eq:C13}
\]

Multiplying by the Gaussian density and integrating gives

\[
 O_{P,J}(h^{2J+2}(1+r)^{3J+3})
 =O_{P,J}((1+r)^{4J+4}/n^{J+1}).
 \tag{C14}\label{tcl:eq:C14}
\]

The real errors in (C9), and the Stirling error, fit inside (C14);
exponentiation preserves those bounds because they tend uniformly to
zero on this window. The convolution remainder (C7) also fits, using
\(H_n/S_n=1+o(1)\).

For clarity about the lattice: \(Z_m=\sqrt{1+r}\,h(m-s)\), so its mesh
is \(\Delta=\sqrt{1+r}\,h\). The normalized central sum is

\[
 \frac{\Delta}{\sqrt{2\pi}}\sum_m e^{-Z_m^2/2}
 e^{E_J(h,Z_m/\sqrt{1+r},r)}Q_J(h,Z_m/\sqrt{1+r},r).
 \tag{C15}\label{tcl:eq:C15}
\]

The mesh-weighted sum of any fixed polynomial times \(e^{-Z^2/2}\) is
uniformly bounded. Thus (C13) controls the summed remainder too. For
each of the finitely many retained polynomial coefficients, Poisson
summation of a Gaussian and its derivatives replaces (C15) by its
integral with error bounded by a polynomial in \(r\) and \(\Delta^{-1}\)
times \(\exp(-2\pi^2/\Delta^2)\), hence smaller than every fixed
required power of \(n\). The actual sum outside \(|Z|\le r^2\) has
relative size \(\exp(-cr^4)\), by the concavity estimates and
\(\tau-s=O(r)\); fixed polynomial factors do not change its
negligibility. This also justifies extending the polynomial-Gaussian
integrals to the whole line.

Finally the analytic expression is invariant under
\((h,y)\mapsto(-h,-y)\). Its coefficient of \(h^q\) has parity \(q\) in
\(y\). Odd terms integrate to zero. The even terms are precisely (C6).
Formula (C12), integrated at \(q=2j\) and multiplied by \((r/2)^j\),
gives \(a_{P,j}=O((1+r)^{4j})\). Gaussian moments put the coefficients
in \(K_P(r)\), and in \(\mathbb{Q}(r)\) for rational \(P\). This proves
(C2).

\subsection{Explicit coefficients and the nonconstant remainder
issue}\label{tcl:explicit-coefficients-and-the-nonconstant-remainder-issue}

The \(S_n\)-normalized first coefficients are

\[
 a_{H,1}=\frac{r^2(r^5-12r^3-28r^2-28r-16)}{48(r+1)^3},
\]

\[
 a_{V,1}=\frac{r^2(r^5-24r^3-64r^2-64r-28)}{48(r+1)^3},
\]

\[
 a_{U,1}=a_{L,1}
 =\frac{r^2(r^5+8r^2+8r-4)}{48(r+1)^3}.
\]

For compact display, define the numerator polynomials
\begin{align*}
N_H(r)={}&r^{11}-48r^9-80r^8+544r^7+2032r^6+2176r^5\\
 &-544r^4-2928r^3-2368r^2-1088r-768,\\
N_V(r)={}&r^{11}-72r^9-152r^8+904r^7+3832r^6+3280r^5\\
 &-7936r^4-22272r^3-23296r^2-12080r-3072,\\
N_U(r)={}&r^{11}-24r^9-8r^8+472r^7+1960r^6+5392r^5\\
 &+12608r^4+20736r^3+20288r^2+10192r+1536.
\end{align*}
The second coefficients are
\[
a_{H,2}=\frac{r^3N_H(r)}{4608(r+1)^6},\qquad
a_{V,2}=\frac{r^3N_V(r)}{4608(r+1)^6},\qquad
a_{U,2}=a_{L,2}=\frac{r^3N_U(r)}{4608(r+1)^6}.
\]

The equality of \(U\) and \(L\) at these two orders is structural:
\(e^{P_L}\) and \(e^{P_U}\) agree through degree two. At the next order,

\[
 a_{L,3}-a_{U,3}=-r^6/48
 \tag{C16}\label{tcl:eq:C16}
\]

exactly. The complete order-three rational functions are in
\texttt{data/coefficients\_order3.json}.

Now define the \textbf{exact} Bell normalization

\[
 T_n=2^{-n}B_{2n}e^{-r^2/4-r/2}.
\]

Then

\[
 \frac{V_n}{T_n}=1+\frac{v_1(r)}n+\frac{v_2(r)}{n^2}
 +\frac{v_3(r)}{n^3}+O(r^{16}/n^4),
\]

where

\[
 v_1(r)=r^2(r-6)(r+3)/48,
\]

\[
\begin{split}
 v_2(r)&=\frac{r^3 N_B(r)}{4608(r+1)^3},\\
 N_B(r)&=r^8-3r^7-66r^6+58r^5+945r^4+813r^3\\
 &\hspace{1cm}-2640r^2-4956r-2208.
\end{split}
\]

and

\[
 v_3(r)\sim r^{12}/663552.
\]

Consequently the remainder after \(v_2/n^2\) is asymptotic to
\(r^{12}/(663552n^3)\). It is \textbf{not \(O(n^{-3})\) with a constant
independent of \(n\)}. A valid simple upper bound at that truncation is
\(O(r^{12}/n^3)\).
\begin{writenote}[How slowly $v_3$ reaches its leading term, 7 October 2026]
The statements $v_3(r)\sim r^{12}/663552$ and the resulting asymptotic form of
the remainder are correct (the write's derivation agrees, and $v_3/r^{12}\to
1/663552$), but the leading term takes over only at astronomically large~$n$.
The numerator of $v_3$ has its largest positive root at $r=10.1555\ldots$:
$v_3(W(2n))<0$ for $2190\le n\le130671$ ($6.5111\ldots<r<10.1555\ldots$),
and $v_3/(r^{12}/663552)$ is $-0.0517$, $-0.0065$, $0.0885$, $0.277$ at
$r=8$, $10$, $12$, $16$, reaching $1/2$ only at $r=23.09\ldots$
($n\approx1.2\times10^{11}$) and $0.9$ at $r=99.08\ldots$
($n\approx5\times10^{44}$). Numerically, with the full $v_3$ the remainder
after $v_2/n^2$ is matched to relative accuracy $10^{-3}$--$10^{-5}$ for
$n=10^4,\ldots,10^7$ (Section~\ref{tcl:sec:provenance}). The bound
$O(r^{12}/n^3)$ stated by the source is the safe reading.
[Independent check, 7~October 2026: the interval above is the last one on
which $v_3<0$. The numerator has five positive roots, all simple, $1.5078\ldots$,
$3.1408\ldots$, $4.6247\ldots$, $6.5111\ldots$ and $10.1555\ldots$. So
$v_3(W(2n))<0$ also for $n\le3$ and for $37\le n\le235$. The sign of the
remainder after $v_2/n^2$ computed from the positive sums agrees: negative at
$n=10^4$ and $10^5$, positive at $10^6$ and $10^7$.]
\end{writenote}


The first relative effect of identifying line images is consequently
\[
\frac{L_n}{U_n}=1-\frac{r^6}{48n^3}+O\left(\frac{r^{16}}{n^4}\right).
\]
This follows by dividing the two expansions with their common first and
second coefficients. The correction separating the legacy series from
corrected line graphs first enters at order four, consistently with the
missing quartic term of the legacy multiplier.
